\chapter{Examples of the Three Information Lifetimes}
\label{app:state-examples}

These examples use the same typed state interface for ordinary Rustlets and
Rustlet Security Domains. Persistent and Session are independent Serde/Postcard
objects. Working data is local to processing and has no inter-APDU retention
contract. None of these categories alone specifies secure memory erasure;
Section~\ref{subsec:rustlet-state-erasure} defines that distinction.

\section{An Ordinary Rustlet with Two Counters}

The following source defines a complete minimal handler. INS~\texttt{10}
increments a durable counter and a session counter; INS~\texttt{20} returns
both as little-endian words. A SELECT does not increment either counter.

\begin{verbatim}
#![no_std]
#![no_main]

use rustlet_runtime::{
    declare_rustlet, Apdu, ApduStatus, Rustlet, RustletCtx, State,
};
use rustlet_runtime::serde::{Deserialize, Serialize};

#[derive(Default, Serialize, Deserialize)]
#[serde(crate = "rustlet_runtime::serde")]
struct Persistent {
    accepted_commands: u32,
}

#[derive(Default, Serialize, Deserialize)]
#[serde(crate = "rustlet_runtime::serde")]
struct Session {
    commands_this_activation: u32,
}

struct Counter;
declare_rustlet!(Counter, 2048usize);

impl Rustlet for Counter {
    type Persistent = Persistent;
    type Session = Session;

    fn process_apdu(
        state: &mut State<Persistent, Session>,
        ctx: &mut RustletCtx,
    ) -> ApduStatus {
        let apdu = Apdu::new(ctx);
        if apdu.is_select() {
            return apdu.accept();
        }
        match apdu.ins() {
            0x10 => {
                state.persistent.accepted_commands = state
                    .persistent.accepted_commands.wrapping_add(1);
                state.session.commands_this_activation = state
                    .session.commands_this_activation.wrapping_add(1);
                apdu.accept()
            }
            0x20 => {
                // Working data: no reference escapes this invocation.
                let mut reply = [0u8; 8];
                reply[..4].copy_from_slice(
                    &state.persistent.accepted_commands.to_le_bytes());
                reply[4..].copy_from_slice(
                    &state.session.commands_this_activation.to_le_bytes());
                apdu.as_sending().send(&reply)
            }
            _ => apdu.reject(ApduStatus::instruction_not_supported()),
        }
    }
}
\end{verbatim}

Only \texttt{accepted\_commands} is encoded in the registry. Successful APDU
publication makes its update durable. The session counter survives later
APDUs and repeated selection of this instance. Selecting another ordinary
instance and returning starts a fresh session counter while preserving the
committed durable counter. Reset has the same separation once Persistent is
recovered from storage. The response array is working data: the transport
owns any staged response bytes after the handler returns, not this local
array or a borrow into it.

The handler type \texttt{Counter} has no instance fields and does not derive
Serde or Default. The runtime owns the two data objects. An application with
no durable state may use \texttt{type Persistent = ();}; one with no session
uses \texttt{type Session = ();}. The typed interface remains the same.

\section{A Security Domain Using the Same State Interface}

This independent source is a minimal administrative skeleton, not an SCP
implementation. It has a durable diagnostic counter and an activation-local
counter. Default secure-channel hooks do not establish a channel. The reset
hook illustrates logical session reset using non-secret diagnostic values;
a real channel also needs deliberate cleanup of its secret storage.

\begin{verbatim}
#![no_std]
#![no_main]

use rustlet_runtime::{
    declare_security_domain, Apdu, ApduStatus, Rustlet,
    RustletCtx, RustletSecurityDomain, State,
};
use rustlet_runtime::serde::{Deserialize, Serialize};

#[derive(Default, Serialize, Deserialize)]
#[serde(crate = "rustlet_runtime::serde")]
struct Persistent {
    accepted_commands: u32,
}

#[derive(Default, Serialize, Deserialize)]
#[serde(crate = "rustlet_runtime::serde")]
struct Session {
    commands_since_reset: u32,
}

struct DiagnosticSd;
declare_security_domain!(DiagnosticSd, 2048usize);

impl Rustlet for DiagnosticSd {
    type Persistent = Persistent;
    type Session = Session;

    fn process_apdu(
        state: &mut State<Persistent, Session>,
        ctx: &mut RustletCtx,
    ) -> ApduStatus {
        let apdu = Apdu::new(ctx);
        if apdu.is_select() {
            return apdu.accept();
        }
        if apdu.ins() != 0x10 {
            return apdu.reject(ApduStatus::instruction_not_supported());
        }
        state.persistent.accepted_commands = state
            .persistent.accepted_commands.wrapping_add(1);
        state.session.commands_since_reset = state
            .session.commands_since_reset.wrapping_add(1);
        apdu.accept()
    }
}

impl RustletSecurityDomain for DiagnosticSd {
    fn reset_secure_channel(state: &mut State<Persistent, Session>) {
        state.session = Session::default();
    }
}
\end{verbatim}

A complete SCP implementation puts derived keys, authentication flags,
handshake progress, MAC chains and encryption counters in Session. Durable
policy stays in Persistent; runtime-owned privileges and lifecycle remain
in the administrative header. Cryptographic scratch values and decoded APDU
borrows remain working data. Protocol handlers must not store APDU borrows
or crypto-provider references in either retained object.

The default Session must represent an inactive channel. A protocol reset
must invalidate channel state immediately even when the SD remains resident.
The assignment shown above resets logical diagnostic values; it is not an
example of guaranteed secret zeroization.

\section{Reading a Rollback Trace}

This is the Persistent-rollback/Session-conserved cell of
Section~\ref{subsec:registry-state-transitions}. A channel reset without a
persistent mutation selects the opposite cell; a terminal SD fault selects
both rollback and session retirement; normal successful processing selects
neither. Session retirement never restores an old authenticated state.


Suppose committed Persistent contains a policy value $P_0$, while the live
Session contains a valid MAC chain $S_0$:
\begin{enumerate}
\item A valid protected command advances the channel to $S_1$.
\item Management tentatively changes the policy to $P_1$, then rejects the
  operation and rolls back its savepoint.
\item The next SD call must observe $P_0$ and the legitimate current Session,
  including $S_1$ and any subsequent valid response-protection progress.
\item The runtime preserves that current Session in a temporary RAM encoding,
  cleans the heap and reconstructs the two typed objects before the call.
  It does not put $S_1$ in flash and does not restore $S_0$.
\end{enumerate}

An unrelated scratch value has no reconstruction guarantee. It may be
removed by heap cleaning or may leave bytes elsewhere in RAM; neither
possibility is application state. A snapshot/restore failure, an SD terminal
fault or final publication failure instead retires the channel. Recovery
requires a fresh activation or establishment as appropriate.

The same distinction applies to a multi-command candidate value. Put it in
Session if losing preparation when the activation ends is acceptable. Put it
in Persistent together with its phase and recovery information if preparation
must survive reset. A local variable or an excluded serialized field cannot
substitute for either contract.
